\chapter{Discussion}
\label{ch:discussion}

This chapter interprets the three substantive findings from Chapters~\ref{ch:replication} through~\ref{ch:ablations} and positions each within the related-work threads from Chapter~\ref{ch:background}. The discussion is organised around three claims: that the persona-vector replication adds a useful (backbone, quantisation) cell to the literature, that the multi-turn drift-detection refutation is the project's most actionable empirical contribution, and that the mechanism cluster surfaces a capability-bound regime that subsequent work on retrieval-side persona conditioning needs to engage with explicitly.

\section{The persona-vector replication in context}

Experiment 1 sits inside the trajectory of recent representation-engineering work that finds linear directions for high-level constructs in LLM hidden states. Marks and Tegmark's geometry-of-truth work\autocite{marks2023geometry} and Arditi's refusal-direction\autocite{arditi2024refusal} are the cleanest published examples of the pattern. The Anthropic persona-vectors paper is the persona-specific instance the project replicates\autocite{chen2025personavectors}, and the Assistant Axis line extends the same family of probes to production-scale models\autocite{lu2026assistantaxis}.

The contribution of the replication is the specific (backbone, quantisation) cell: Gemma-2-9B-Instruct at 4-bit NF4 on V100. The Anthropic paper validates on Qwen 2.5-7B and Llama-3.1-8B in fp16. The Assistant Axis work runs at larger scales but does not test 4-bit quantisation systematically. The replication closes one cell of a sweep that subsequent work needs anyway, with the Dubanowska controls intact and the contrast-prompt set, the cache schema, and the validation script all available in the project repository for direct reproduction\autocite{dubanowska2025}.

The honest framing on the replication's strength is that the test set is small ($n = 10$ per persona on the prompt-disjoint split) and the \auroc{} 1.000 saturation is partly an artefact of the contrast-prompt formulation rather than a strong claim about layer-by-layer separability. The within-prompt explicit role instruction is recoverable linearly at every layer in the swept band; Experiment 1 confirms this and Experiment 2 demonstrates that the same recoverability does not extend to multi-turn dialogue without an explicit role instruction in the prompt prefix. The two experiments together are sharper than either alone.

\section{What the multi-turn drift-detection refutation means}

Experiment 2 is the project's strongest single empirical finding. The result is a refutation of a specific architectural assumption: that the contrast-prompt-trained persona direction transfers to inference-time drift detection in multi-turn dialogue at quantised inference scale on the production-scale instruct-tuned models the project tested. The result is not a refutation of the published persona-vectors methodology. Experiment 1 replicates cleanly. The direction is real. What does not hold is the unstated assumption, implicit in the v0.3 PRD design and absent from the persona-vectors paper's validation, that the same direction would also serve as a drift-detection signal at inference time.

\subsection{Why the result is robust}

The empirical coverage is the contribution. 24 generation-scope cells (2 backbones $\times$ 4 layers $\times$ 3 personas in generation-scope) plus 12 prompt-scope cells on Gemma, every cell refuted by the pre-registered $+0.30$ threshold. The maximum signal observed is 0.103, an order of magnitude below threshold and in the wrong direction. Roughly half the cells show the signal flipped relative to the expected direction, which is consistent with the direction being approximately orthogonal to the construct we cared about rather than aligned-but-weak.

Three hypotheses about the cause remain open, as Chapter~\ref{ch:replication} flagged. The contrast-prompt setup may encode role-instruction-presence rather than generation-content-fidelity. 4-bit quantisation may disrupt the direction's usefulness in ways it does not disrupt contrast-prompt separability. The prompt-versus-generation distinction may be fundamentally geometric in these models. Distinguishing the three is thesis-bridge work and would, ideally, be answered before another round of drift-gated mechanism design at this scale. The cleanest next experiment is a generation-scope re-extraction trained on hidden states sampled from genuinely drifting dialogue rather than from contrast prompts. The drift-trajectory corpus this project authored is a natural starting point.

\subsection{Why this matters for retrieval-side persona conditioning}

The architectural pattern Skill-RAG\autocite{wei2026skillrag} and Probing-RAG\autocite{baek2025probingrag} pioneered, gating expensive retrieval-time machinery on an internal-state signal, depends critically on the signal being differentially meaningful at inference time. Experiment 2 establishes that the specific signal the v3 Persona-RAG design proposed, the persona-vector projection, does not have this property at the scale and quantisation the project tested. The implication is not that Skill-RAG's broader pattern is wrong; Skill-RAG's signal (failure-state detection) is different from the persona-vector projection and was independently validated by Wei and colleagues.

The implication is narrower: the persona-vectors literature's validation methodology, while internally clean, does not generalise to multi-turn drift detection as easily as the v0.3 PRD assumed. Subsequent persona-conditioned RAG work that proposes to use persona vectors as a gate signal should test the multi-turn transfer separately, with the kind of empirical sweep Chapter~\ref{ch:replication} describes. Two hand-authored conditions on three personas over six turns, with both extraction regimes and four layers per cell, is around 36 forward passes per (backbone, scope) cell. The cost is small. The architectural saving is large.

\section{The capability-bound mechanism cluster}

The mechanism cluster reported in Chapter~\ref{ch:mechanism-results} and probed in Chapter~\ref{ch:ablations} is the project's third substantive finding. B2, M1, and M3 cluster within 0.20 PoLL points on every probe type at Gemma-2-9B-Instruct. The three diagnostic ablations isolate three plausible confounds (gate calibration, ID-RAG, quantisation precision) and none of them moves the cluster outside the pre-registered envelope. The convergent reading is that the binding constraint at this scale is model capability, not retrieval architecture.

\subsection{What the cluster says about retrieval architectures}

The defensible claim is that on this benchmark, at this scale, with the specific metric stack the project used, the architectural headroom above a properly-engineered RoleGPT-level prompt baseline is small. The honest framing is that this does not invalidate retrieval-side persona conditioning; it bounds what one can credibly claim from a 9B-scale evaluation. Two observations support the narrower interpretation.

First, the B1 collapse is robust. PoLL persona-adherence at 2.54 on B1 against the cluster's 4.5 to 4.7 is a $+2$-point gap on a 1-to-5 scale, type-independent and persona-independent. The persona-conditioning project's premise survives without difficulty: persona has to be encoded somewhere, the question is just where. The cluster says only that on this benchmark, encoding the persona in the prompt structure and in the typed-memory schema produce indistinguishable outputs at the PoLL panel's scoring granularity.

Second, the Type A gain at fp16 (Section~\ref{sec:quant-sweep}, $-0.184$ $\Delta$-of-$\Delta$, M1 gains over B2) is suggestive but small. A larger-scale evaluation (3 personas $\times$ 30 conversations rather than 1 persona $\times$ 10 conversations) would test whether the gain is real and whether it extends past the self-fact-challenge probe type. The thesis-bridge follow-up is named in Chapter~\ref{ch:limitations}.

\subsection{Why M3's gate fires conservatively}

The M3 gate's precision-recall asymmetry (P=1.00 / R=0.06 on the probe corpus at the pilot-calibrated threshold) is a calibration finding, not an architectural one. The 0.5 confidence threshold inherited from the drift-trajectory pilot was calibrated against conversations that triggered the gate at around 40\%. The probe corpus produces different judge-confidence distributions because the turns are content-shorter, more single-topic, and more probe-shaped than the drift-trajectory conversations. The judge's confidence stays high on most probe turns and clusters low on the small fraction where the user's reframing or the injected counter-evidence is unambiguous. The threshold-sensitivity sweep at $\{0.2, 0.3, 0.4, 0.5\}$ on a 5-probe sample is the cleanest follow-up to characterise the precision-recall tradeoff and to identify whether a corpus-tuned threshold makes the cost premium architecturally justified. The sweep is named as thesis-bridge follow-up \#1.

\section{What the two-experiment design contributes methodologically}

The two-experiment design works in this report partly because the experiments answer adjacent questions and partly because the second experiment refutes an assumption the first experiment alone could not have tested. Reporting either experiment without the other would have produced a substantively different paper.

A persona-vector paper that reported only Experiment 1's replication on a new (backbone, quantisation) cell would be a respectable but narrow methods note. The contribution would be a single (backbone, quantisation) data point that subsequent work would mostly cite as confirming a known recipe. A persona-vector paper that reported only Experiment 2's refutation, without Experiment 1's clean replication, would be much harder to interpret. The reader would suspect the methodology was implemented incorrectly, or that the contrast prompts were too weak, or that the layer choice was wrong. The clean Experiment 1 result is what makes Experiment 2's refutation the interesting finding it is. The Dubanowska controls in Experiment 1 are the second layer of methodological discipline that gives Experiment 2 its weight, because a reader can verify that the direction is not a spurious correlation before the multi-turn transfer is tested.

The methodological lesson generalises. When the published validation of an inference-time signal uses one regime (contrast prompts, single-turn) and the proposed application uses another regime (multi-turn dialogue), the transfer is a real empirical question and should be tested directly. The cost is small. The architectural saving, if the transfer fails, is large.

\section{Limitations of the findings themselves}

Three limitations qualify the chapter's main claims. The PoLL panel's panel-versus-human Krippendorff $\alpha$ of 0.306 (Chapter~\ref{ch:methodology}) is YELLOW per the pre-registered thresholds. The mechanism cluster's interpretation depends on the panel's relative ordering being faithful to what a human-evaluator panel would have produced, which the 20-item pilot does not strongly support. A 150-item human-validation expansion is named as future work. The MiniCheck false-positive issue (Section~\ref{sec:minicheck-fp-audit}) limits the interpretability of MiniCheck on M1 and M3 specifically; the headline persona-quality column is PoLL alone. The Spec-9 ablation cell sizes are small ($n = 10$ on the precision sweep specifically); a Type A signal at the suggestive-but-not-strong band ($-0.184$ $\Delta$-of-$\Delta$) is what one would expect from underpowered cells if the effect were real, but it is also what one would expect from underpowered cells if the effect were a sampling artefact. Disentangling the two requires scale-up.

Chapter~\ref{ch:limitations} expands on these and names the follow-up directions that fall out of each.
